(1) 两束反射光的光程差为 \(2t(x)\)；下方空气向玻璃反射附加相位 \(\pi \)，上方玻璃向空气反射无该项。因此总相位差为
\[
\delta(x)=\frac{4\pi t(x)}{\lambda}+\pi
=\frac{4\pi\alpha x}{\lambda}+\pi.
\]
忽略弱反射的幅值差，干涉强度与 \(1+\cos\delta(x)\) 成正比。

(2) 暗纹要求 \(\delta=(2j+1)\pi\)，所以 \(2t=j\lambda\)，
\[
x_j^{\rm dark}=\frac{j\lambda}{2\alpha},\quad j=0,1,\dots.
\]
亮纹要求 \(\delta=2j\pi\)，即 \(2t=(j+1/2)\lambda\)，故
\(x_j^{\rm bright}=(j+1/2)\lambda/(2\alpha)\)。

(3) 相邻暗纹间距
\[
\Delta x=x_{j+1}-x_j=\frac{\lambda}{2\alpha}.
\]
在小角近似下间距与 x 无关，反射条纹形成均匀间隔的等厚干涉带。

(4) 当 x=0 两板理想接触 t=0，几何光程差为0，但反射相位差为 \(\pi \)，故反射暗纹。测得 \(\Delta x\) 可以求 \(\alpha\simeq\lambda/(2\Delta x)\)。若玻璃真实接触处存在污渍或纳米间隙，首纹位置可能偏移，需校准零厚度。
